\documentclass[11pt,letterpaper]{article}

\usepackage[T1]{fontenc}
\usepackage[utf8]{inputenc}
\usepackage{charter}
\usepackage[scaled=0.88]{helvet}
\usepackage{courier}
\usepackage{microtype}
\usepackage{amsmath,amssymb}
\usepackage{booktabs}
\usepackage{longtable}
\usepackage{array}
\usepackage{enumitem}
\usepackage{listings}
\usepackage{tcolorbox}
\usepackage{fancyhdr}
\usepackage{needspace}
\usepackage{titlesec}
\usepackage[table]{xcolor}
\usepackage[hidelinks]{hyperref}
\tcbuselibrary{skins,breakable}

\usepackage[letterpaper,top=1.0in,bottom=1.0in,left=1.1in,right=1.1in,
            headheight=15pt,footskip=0.5in]{geometry}

% ------------------------------------------------------------------ colour
\definecolor{accent}{HTML}{EB6834}
\definecolor{cool}{HTML}{2A78D6}
\definecolor{ink}{HTML}{1B1B1B}
\definecolor{muted}{HTML}{898781}
\definecolor{good}{HTML}{1C7C54}
\definecolor{rule}{HTML}{D8D6D1}

% ------------------------------------------------------------------ headings
\titleformat{\section}{\sffamily\large\bfseries\color{ink}}{\thesection.}{0.6em}{}
\titleformat{\subsection}{\sffamily\small\bfseries\color{muted}}{}{0em}{}
\titlespacing*{\section}{0pt}{22pt plus 4pt minus 2pt}{8pt}
\titlespacing*{\subsection}{0pt}{12pt plus 3pt minus 2pt}{4pt}

\setlist[itemize]{itemsep=3pt,topsep=4pt,parsep=0pt,leftmargin=1.4em}
\setlist[enumerate]{itemsep=3pt,topsep=4pt,parsep=0pt,leftmargin=1.6em}
\setlength{\parskip}{5pt}
\setlength{\parindent}{0pt}

% ------------------------------------------------------------------ header
\pagestyle{fancy}
\fancyhf{}
\renewcommand{\headrulewidth}{0.4pt}
\renewcommand{\headrule}{\color{rule}\hrule height 0.4pt}
\fancyhead[L]{\sffamily\scriptsize\color{muted}MATH 153 \textemdash{} Data Wrangling / Data Engineering}
\fancyhead[R]{\sffamily\scriptsize\color{muted}Building a Simple ETL Pipeline}
\fancyfoot[C]{\sffamily\scriptsize\color{muted}\thepage}

% ------------------------------------------------------------------ code
\lstdefinestyle{py}{
  language=Python,
  basicstyle=\ttfamily\small,
  keywordstyle=\color{accent},
  stringstyle=\color{good},
  commentstyle=\color{muted}\itshape,
  showstringspaces=false,
  breaklines=true,
  frame=leftline,
  framesep=7pt,
  rulecolor=\color{rule},
  aboveskip=8pt,
  belowskip=6pt,
  xleftmargin=10pt,
}
\lstset{style=py}
\lstdefinestyle{sql}{language=SQL,morekeywords={NOT,EXISTS,IN,LEFT,JOIN,HAVING,PRAGMA}}

\newcommand{\py}[1]{{\ttfamily\small#1}}

% ------------------------------------------------------------------ boxes
\newtcolorbox{ask}[1][]{
  enhanced, breakable, colback=accent!5, colframe=accent,
  boxrule=0pt, leftrule=3pt, arc=1pt,
  left=10pt, right=10pt, top=7pt, bottom=7pt,
  fonttitle=\sffamily\bfseries\footnotesize, #1}

\newtcolorbox{trap}[1][]{
  enhanced, breakable, colback=muted!7, colframe=muted,
  boxrule=0pt, leftrule=2.5pt, arc=1pt,
  left=9pt, right=9pt, top=6pt, bottom=6pt,
  fonttitle=\sffamily\bfseries\footnotesize, title={Watch out}, #1}

\newtcolorbox{lesson}[1][]{
  enhanced, breakable, colback=good!5, colframe=good,
  boxrule=0pt, leftrule=2.5pt, arc=1pt,
  left=9pt, right=9pt, top=6pt, bottom=6pt,
  fonttitle=\sffamily\bfseries\footnotesize, title={Remember}, #1}

\newtcolorbox{keynote}[1][]{
  enhanced, breakable, colback=cool!4, colframe=cool,
  boxrule=0pt, leftrule=2.5pt, arc=1pt,
  left=9pt, right=9pt, top=6pt, bottom=6pt, #1}

\newcommand{\tests}[1]{{\sffamily\footnotesize\color{muted}\textbf{Probes:} #1}\par\vspace{2pt}}
\newcommand{\ans}{\subsection{The answer}}

% --- extra boxes for this handout
\newtcolorbox{greenbox}[1][]{
  enhanced, breakable, colback=good!5, colframe=good,
  boxrule=0pt, leftrule=2.5pt, arc=1pt,
  left=9pt, right=9pt, top=6pt, bottom=6pt,
  fonttitle=\sffamily\bfseries\footnotesize, title={Remember}, #1}

\newtcolorbox{greybox}[1][]{
  enhanced, breakable, colback=muted!7, colframe=muted,
  boxrule=0pt, leftrule=2.5pt, arc=1pt,
  left=9pt, right=9pt, top=6pt, bottom=6pt, #1}
\begin{document}

\begin{center}
  {\sffamily\LARGE\bfseries\color{ink} Building a Simple ETL Pipeline}\\[4pt]
  {\sffamily\large\color{accent} Components of Data Pipelines}\\[9pt]
  {\sffamily\small\color{muted} MATH 153 \textbar{} Data Wrangling / Data Engineering
   \textbar{} Howard University, Department of Mathematics}\\[2pt]
  {\sffamily\small\color{muted} Moussa Doumbia, Ph.D.}
\end{center}

\vspace{6pt}
\color{rule}\hrule height 0.4pt\color{ink}
\vspace{8pt}

{\sffamily\small\color{muted}\textbf{Files for this lesson}}

{\sffamily\footnotesize
\begin{tabular}{@{}p{6.1cm}p{7.8cm}@{}}
\py{ETL\_Pipeline\_Lecture.ipynb} & the worked example, ready to run \\
\py{ETL\_Activity.ipynb}         & the in-class activity (with blanks to fill in) \\
\py{ETL\_Activity\_SOLUTIONS.ipynb} & the answer key \\
\py{data/}                       & caf\'e sales: 3 daily files + \py{menu.csv} \\
\py{activity\_data/}             & library checkouts: 3 daily files + \py{books.csv} \\
\end{tabular}}

% ==================================================================
\section{What is a data pipeline?}

\begin{keynote}
A \textbf{data pipeline} is a series of steps that moves data from where it \emph{is} to where it is
\emph{useful}, doing the same thing every time it runs.
\end{keynote}

The important words are \textbf{every time it runs}. A notebook you run once by hand is an analysis.
A pipeline runs tomorrow, and next Tuesday, without you.

\subsection{The five components}

{\sffamily\footnotesize
\begin{tabular}{@{}p{0.5cm}p{2.3cm}p{5.0cm}p{5.4cm}@{}}
\toprule
\textbf{\#} & \textbf{Component} & \textbf{What it means} & \textbf{In our example} \\
\midrule
1 & \textbf{Source}    & where the data comes from       & 3 daily CSV files from the caf\'e till \\
2 & \textbf{Extract}   & read it in                      & \py{pd.read\_csv} \\
3 & \textbf{Transform} & make it usable                  & fix prices, fix names, join the menu \\
4 & \textbf{Load}      & put it where people will use it & a SQLite table + a summary CSV \\
5 & \textbf{Schedule}  & decide when it runs             & every morning at 6 a.m. \\
\bottomrule
\end{tabular}}

\vspace{6pt}
Steps 2, 3 and 4 are so common that they have a name of their own: \textbf{ETL}.

\begin{lstlisting}[language=]
 SOURCE      EXTRACT     TRANSFORM      LOAD       SCHEDULE
 3 CSVs  --> read    --> clean+join --> cafe.db <-- daily
 + menu                                 summary     6 a.m.
\end{lstlisting}

\subsection{ETL in one sentence each}

\begin{itemize}
  \item \textbf{Extract} \textemdash{} get the raw data, \emph{without changing it}.
  \item \textbf{Transform} \textemdash{} clean it, reshape it, combine it, calculate from it.
  \item \textbf{Load} \textemdash{} write the result somewhere other people can read.
\end{itemize}

\begin{greybox}
\small
You will also hear \textbf{ELT}: load the raw data into the database first, and transform it later
\emph{inside} the database. Same three letters, different order. ELT suits very large data and
powerful warehouses; ETL suits what we are doing today.
\end{greybox}

% ==================================================================
\section{The worked example: Bison Caf\'e sales}

The caf\'e sends one sales file per day. We have three days. We want a clean table and a daily
revenue report.

\subsection{Look at the data first}

\emph{Never write a pipeline for data you have not looked at.}

\begin{lstlisting}[language=]
sale_id,sale_date,sale_time,item,qty,unit_price
S-5001,2027-03-02,14:32,FRUIT CUP,2,$4.00
S-5002,2027-03-02,15:22,Turkey Wrap,2,$7.50
S-5004,2027-03-02,11:58,latte,3,$4.25
\end{lstlisting}

Two things are already wrong, and both are normal:

\begin{enumerate}
  \item \py{unit\_price} \textbf{is text, not a number}. It says \py{\$4.25}. You cannot multiply by
        \py{\$4.25}.
  \item \py{item} \textbf{is spelled inconsistently} \textemdash{} \py{FRUIT CUP}, \py{latte},
        \py{"  Bagel "} \textemdash{} because the till lets staff type it by hand.
\end{enumerate}

Fixing those two things is what ``Transform'' will mean today. There is also a reference table,
\py{menu.csv}, giving each item's category.

\subsection{E \textemdash{} Extract}

Read the raw files. Change nothing.

\begin{lstlisting}
def extract(data_dir):
    """Read every daily sales file and stack them into one DataFrame."""
    files = sorted(data_dir.glob("sales_2027-*.csv"))
    frames = [pd.read_csv(f) for f in files]
    return pd.concat(frames, ignore_index=True)

raw = extract(DATA)        # 107 rows from 3 files
\end{lstlisting}

\begin{greenbox}[title={Remember}]
\small
\textbf{Extract is deliberately boring.} Keeping the raw data untouched means that when something goes
wrong later, you know where to start again. \py{ignore\_index=True} renumbers the stacked rows
0, 1, 2, \ldots{} so you do not end up with three rows all numbered 0.
\end{greenbox}

\subsection{T \textemdash{} Transform}

Four small steps.

\vspace{4pt}
\textbf{Step 1 \textemdash{} make the price a number.}
\begin{lstlisting}
df["unit_price"] = df["unit_price"].str.replace("$", "", regex=False).astype(float)
\end{lstlisting}

\textbf{Step 2 \textemdash{} make the date a real date.}
\begin{lstlisting}
df["sale_date"] = pd.to_datetime(df["sale_date"], format="%Y-%m-%d")
\end{lstlisting}

A date stored as text is just a string \textemdash{} you cannot sort it reliably, subtract it, or ask
what day of the week it was.

\textbf{Step 3 \textemdash{} clean the item names.}
\begin{lstlisting}
df["item"] = df["item"].str.strip().str.title()
\end{lstlisting}

\begin{lstlisting}[language=]
distinct spellings before: 29
distinct item names after:  8
\end{lstlisting}

\py{.str.strip()} removes spaces at the ends; \py{.str.title()} capitalises each word, so
\py{"LATTE"} and \py{"latte"} both become \py{"Latte"}.

\textbf{Step 4 \textemdash{} join the menu, then calculate.}
\begin{lstlisting}
df = df.merge(menu, on="item", how="left")
df["revenue"] = (df["qty"] * df["unit_price"]).round(2)
\end{lstlisting}

\begin{trap}
\small
\textbf{Step 3 had to come before step 4.} \py{"latte"} would not have matched \py{"Latte"} in the
menu, that sale would have had no category, and its revenue would have quietly vanished from the
report.
\end{trap}

\begin{greenbox}[title={Two checks to make after every join}]
\small
\textbf{Did the row count stay the same?} 107 in, 107 out. If it grew, the reference table has
duplicate rows and every sale got counted twice.

\medskip
\textbf{Did anything fail to match?} Zero unmatched here. A blank category means a spelling your
cleaning step missed.
\end{greenbox}

\subsection{L \textemdash{} Load}

Two outputs, for two audiences: a \textbf{database table} for anyone who wants to query it, and a
\textbf{summary CSV} for the manager who just wants the daily numbers.

\begin{lstlisting}
con = sqlite3.connect("cafe.db")
out.to_sql("sales", con, if_exists="replace", index=False)
con.close()
summary.to_csv("daily_summary.csv", index=False)
\end{lstlisting}

The result, from 107 sales over three days:

{\sffamily\footnotesize
\begin{tabular}{@{}llrr@{}}
\toprule
\textbf{sale\_date} & \textbf{category} & \textbf{items\_sold} & \textbf{revenue} \\
\midrule
2027-03-02 & Drink & 35 & \$119.75 \\
2027-03-02 & Food  & 40 & \$222.85 \\
2027-03-03 & Drink & 48 & \$162.80 \\
2027-03-03 & Food  & 29 & \$112.40 \\
2027-03-04 & Drink & 39 & \$124.30 \\
2027-03-04 & Food  & 29 & \$126.25 \\
\midrule
\multicolumn{3}{@{}l}{\textbf{Total}} & \textbf{\$868.35} \\
\bottomrule
\end{tabular}}

% ==================================================================
\section{The whole pipeline, in one function}

Everything above, in twelve lines. \textbf{This is the deliverable} \textemdash{} not the notebook,
the function.

\begin{lstlisting}
def run_pipeline(data_dir=DATA, db_path="cafe.db", csv_path="daily_summary.csv"):
    """Extract -> Transform -> Load. Run this every morning."""
    menu  = pd.read_csv(data_dir / "menu.csv")

    raw   = extract(data_dir)                       # E
    clean = transform(raw, menu)                    # T

    assert clean["category"].notna().all(), "some items did not match the menu"
    assert (clean["revenue"] > 0).all(),    "a sale has zero or negative revenue"

    summary = (clean
               .assign(day=clean["sale_date"].dt.date)
               .groupby(["day", "category"], as_index=False)
               .agg(items_sold=("qty", "sum"), revenue=("revenue", "sum"))
               .rename(columns={"day": "sale_date"}))

    load(clean, summary, db_path, csv_path)         # L
    return clean, summary
\end{lstlisting}

\begin{greenbox}[title={Why a function, and not just cells?}]
\small
\textbf{It has a name.} You can tell someone ``run \py{run\_pipeline()}'' instead of ``open the
notebook and run cells 4 through 19 but not cell 12''.

\medskip
\textbf{It takes arguments.} Point it at a different folder and it works on different data.

\medskip
\textbf{It can be scheduled.} A computer can call a function at 6 a.m. It cannot click through a
notebook.
\end{greenbox}

% ==================================================================
\section{Two things every pipeline needs}

\subsection{Run it twice}

\begin{lstlisting}[language=]
run_pipeline()
run_pipeline()          # on purpose

rows in the sales table: 107        <-- not 214
\end{lstlisting}

Running it twice gave the same answer as running it once. That property has a name: the pipeline is
\textbf{idempotent}. One keyword gives it to us \textemdash{} \py{if\_exists="replace"} throws away
the old table and writes a fresh one each run.

\begin{trap}
\small
\textbf{The most common bug in a first pipeline.}
\begin{lstlisting}
out.to_sql("sales", con, if_exists="append", index=False)   # DON'T
\end{lstlisting}
With \py{append}, every run adds another copy. Run it three times and the caf\'e has sold everything
three times.
\end{trap}

\subsection{A log line per step}

\begin{lstlisting}[language=]
  EXTRACT    107 rows read
  TRANSFORM  107 rows cleaned, 8 distinct items
  LOAD       107 rows written, revenue $868.35
\end{lstlisting}

Three \py{print} statements. Now when the job runs at 6 a.m. and something looks odd, you can see
\emph{which step} the rows disappeared at instead of guessing.

\subsection{A check before you publish}

An \textbf{assertion} is a sentence that must be true. If it is not, Python stops immediately and
tells you why.

\begin{lstlisting}
assert clean["category"].notna().all(), "some items did not match the menu"
assert (clean["revenue"] > 0).all(),    "a sale has zero or negative revenue"
\end{lstlisting}

\begin{greenbox}[title={Where the checks go}]
\small
\textbf{Between Transform and Load.} That is the last moment before the numbers reach somebody else.
If a check fails, the pipeline stops \emph{before} writing to the database, and the manager's report
is still yesterday's \textemdash{} which is correct, and much better than today's report being wrong.
\end{greenbox}

% ==================================================================
\section{The same thing in SQL}

Once the data is in the database, the summary is one query.

\begin{lstlisting}[style=sql]
SELECT sale_date, category,
       SUM(qty)               AS items_sold,
       ROUND(SUM(revenue), 2) AS revenue
FROM   sales
GROUP  BY sale_date, category
ORDER  BY sale_date, category;
\end{lstlisting}

{\sffamily\footnotesize
\begin{tabular}{@{}p{6.4cm}p{6.4cm}@{}}
\toprule
\textbf{Use pandas when} & \textbf{Use SQL when} \\
\midrule
reading files and cleaning text       & the data is already in a database \\
the step is awkward to express in SQL & the data is too big for memory \\
you are still exploring               & the question is a standard group-and-total \\
\bottomrule
\end{tabular}}

\vspace{4pt}
Most real pipelines use both: pandas to clean and load, SQL to query afterwards. That is exactly what
we did.

% ==================================================================
\section{Five things to remember}

\begin{enumerate}
  \item \textbf{Extract changes nothing.} Keep the raw data raw, so you can always start over.
  \item \textbf{Clean before you join.} \py{"latte"} does not match \py{"Latte"}.
  \item \textbf{Check the row count after every join.} Same number in, same number out.
  \item \textbf{Make it runnable twice.} \py{if\_exists="replace"}, not \py{"append"}.
  \item \textbf{Put the checks between Transform and Load} \textemdash{} the last moment before the
        numbers reach somebody else.
\end{enumerate}

\begin{keynote}
A pipeline is just \textbf{three functions and a schedule}. The difference between a pipeline and a
notebook is that somebody else can run it.
\end{keynote}

% ==================================================================
\newpage
\section{In-class activity: build a basic data pipeline}

{\sffamily\small\color{muted}\textbf{Notebook:} \py{ETL\_Activity.ipynb} \textbullet{}
\textbf{Work in pairs} \textbullet{} \textbf{About 35 minutes}}

\subsection*{\sffamily The situation}

Founders Library records every book checkout and sends the circulation office \textbf{one file per
day}. Three days have arrived. The office wants a clean table they can query, plus a daily summary of
how many books went out by genre.

\textbf{Your job: build the ETL pipeline.} Same five components as the caf\'e example \textemdash{}
only the data is different.

{\sffamily\footnotesize
\begin{tabular}{@{}p{2.2cm}p{11.5cm}@{}}
\toprule
\textbf{Source}    & \py{activity\_data/checkouts\_2027-03-0\{2,3,4\}.csv} and \py{books.csv} \\
\textbf{Extract}   & read the three daily files and stack them \\
\textbf{Transform} & fix \py{loan\_days}, clean \py{patron\_type}, join \py{books}, add a due date \\
\textbf{Load}      & write to \py{library.db} and \py{checkout\_summary.csv} \\
\textbf{Schedule}  & (we will just talk about this one) \\
\bottomrule
\end{tabular}}

\subsection*{\sffamily The tasks}

{\sffamily\footnotesize
\begin{tabular}{@{}p{0.7cm}p{1.1cm}p{5.6cm}p{5.6cm}@{}}
\toprule
\textbf{\#} & \textbf{Time} & \textbf{Task} & \textbf{Expected result} \\
\midrule
1 & 5 min & \textbf{Extract} \textemdash{} write a function that reads all three daily files and
                stacks them & 82 rows \\
2 & 5 min & \textbf{Transform 1} \textemdash{} turn \py{"14 days"} into the number \py{14};
                convert \py{checkout\_date} to a datetime & dtype \py{int64}, not \py{object} \\
3 & 5 min & \textbf{Transform 2} \textemdash{} make \py{patron\_type} consistent
            & \textbf{16 spellings become 4} \\
4 & 7 min & \textbf{Transform 3} \textemdash{} left-join \py{books} on \py{book\_id}; compute
                \py{due\_date = checkout\_date + loan\_days} & still 82 rows, 0 unmatched \\
5 & 6 min & \textbf{Load} \textemdash{} build the daily summary; write to \py{library.db} and a CSV
            & 82 rows loaded \\
6 & 8 min & \textbf{Put it together} \textemdash{} one \py{run\_pipeline()} function, a log line per
                step, and \textbf{two assertions} between Transform and Load & runs twice, still 82
                rows \\
\midrule
\multicolumn{4}{@{}l}{\textit{Bonus:} query the database \textemdash{} which genre was checked out
                      most? How many checkouts per patron type?} \\
\bottomrule
\end{tabular}}

\subsection*{\sffamily What to look for while they work}

\begin{trap}
\small
\textbf{Task 3 before Task 4.} A pair who joins first and cleans second will get the right row count
and the wrong idea. Ask them: \emph{``why did this step have to come before the join?''}
\end{trap}

\begin{trap}
\small
\textbf{Task 6, running it twice.} A pair who used \py{if\_exists="append"} will see
\textbf{246 rows} instead of 82. This is the single most valuable mistake in the activity
\textemdash{} let them find it, then ask what the library director would conclude from a report
saying the collection circulated three times over.
\end{trap}

\begin{trap}
\small
\textbf{\py{pd.to\_timedelta} needs a unit.} For the due date:
\py{pd.to\_timedelta(df["loan\_days"], unit="D")}. Without \py{unit="D"} it will not know whether
you meant days, hours or nanoseconds.
\end{trap}

\subsection*{\sffamily Three questions to close on}

\begin{enumerate}
  \item \textbf{Scheduling.} Your pipeline reads every file in the folder. Tomorrow a fourth file
        arrives. What happens when you run \py{run\_pipeline()} again \textemdash{} and is that what
        you want?
  \item \textbf{Where checks belong.} You put two assertions between Transform and Load. What would
        go wrong if you put them at the very end, after \py{load()} had already run?
  \item \textbf{One thing you would add.} You have a working pipeline. Name one thing it still needs
        before you would be comfortable letting it run unattended every morning.
\end{enumerate}

\vspace{6pt}
\begin{keynote}
\small
\textbf{What they built:}
\py{extract()} \(\rightarrow\) transform \(\rightarrow\) \py{assert} \(\rightarrow\) \py{load()}
\textemdash{} 82 rows, 2 checks, and a result somebody else can use. Four functions and two checks.
That is a data pipeline.
\end{keynote}

\vspace{8pt}
\color{rule}\hrule height 0.4pt\color{ink}
\vspace{4pt}

{\sffamily\footnotesize\color{muted}
Every figure in this document is produced by the accompanying notebooks, which run top to bottom from
a clean kernel. The data are synthetic, generated by \py{make\_data.py} with a fixed seed, and
describe no real caf\'e, library or patron.}

\end{document}
